\section{Theory Development}\label{sec:background}

To theorize about the effects of license choice on a project's performance,
several strands of prior research may be relevant.
First is the line of work on motivations of
individuals~\cite{von2012carrots} and businesses~\cite{Harhoff02} to
participate in OSS.
Second, technical, social, and other factors
that either constraint certain choices in software development or
make them more convenient~\cite{ma2020methodology}.
Third, the specific plans and objectives of the project, and how each license choice may advance and hinder these goals~\cite{sen2008determinants}.

\subsection{Motivations}

\citet{von2012carrots} categorize OSS participant motivations
\icsetrim{into ideology, altruism, kinship, fun, reputation,
reciprocity, learning, own-use, career, and pay}
into ten classes, several of which plausibly influence the type of
licenses contributors prefer (ideology and reciprocity favor strong
copyleft; altruism leans permissive; kinship drives adoption of
licenses already used by collaborators).
\icsetrim{In case of commercial involvement, motivations may diverge
from individual contributors, even if they generally align with profit
maximization~\cite{Harhoff02}. Studies of open and gated software
communities highlight the difficulties firms face in balancing
community-based value creation with private value
appropriation~\cite{Shah06}. However, community sponsorship and
licensing strategies can address these tensions. For instance,
restrictive licenses can ensure long-term control over contributions,
whereas nonmarket\footnote{The term ``nonmarket'' excludes for-profit
organizations.} sponsors might alleviate concerns about the project's
sustainability without the need for restrictive
licenses~\cite{stewart2006impacts}.}
For commercial actors, motivations align with profit maximization
and shape licensing strategy in characteristic
ways~\cite{Harhoff02,Shah06,stewart2006impacts}.
\citet{fershtman2007open} observed empirically that restrictive
licenses attract more contributors per project, while permissive
licenses yield higher productivity per contributor.
\icsetrim{Similarly, research into firm involvement in OSS shows that
licensing choices may be driven by strategic goals. For instance, firms
may adopt permissive licenses to facilitate bundling proprietary and
open source code~\cite{economicos02}. \citet{dahlander2008firms}
identified three strategies firms use to engage with OSS communities:
(1) accessing community development to extend resources, (2) aligning
their strategy with community efforts, and (3) assimilating the
community to integrate and share outcomes. Different commercial models
necessitate distinct licensing strategies. \citet{wagstrom2010impact}
and \citet{zhou2016inflow} identified two types of firm engagement:
community-focused firms (e.g., GNOME), which prioritize building
vibrant OSS communities and monetize through services (which favor
restrictive licenses so that a competitor can not create proprietary
enhancements without sharing the code), and product-focused firms
(e.g., Eclipse), which rely on product revenues and thus favor weaker
copyleft provisions to allow some forms of bundling with proprietary
code.}
Firms differentiate further by engagement model
(\citet{dahlander2008firms}, \citet{wagstrom2010impact},
\citet{zhou2016inflow}): community-focused firms favor restrictive
licenses to prevent third-party proprietary enhancements,
product-focused firms prefer weak copyleft to allow proprietary
bundling.
\icsetrim{Emergence of cloud services endangered the business model of
OSS companies that rely on service contracts. Because cloud companies
provide computing services, they can also run (and sell) the same OSS
software as a service, thus cutting out companies that developed the
software. For example, in 2021, Elasticsearch changed its Apache
license to a commercial one where the products built from that code
can not be provided to others as a managed service. They made another
change in 2024, moving to Affero GPL, which is a copyleft license
that requires anyone running a modified program on a server and
letting other users communicate with it to also provide access to the
source code of the modified version running on the server. As we can
see, the more permissive license (Apache) was changed to a more
restrictive one (commercial, then AGPL), as is typical for the service
model.}

Following \citet{fershtman2007open}'s conjecture:

\begin{itemize}
    \item \textbf{Hypothesis (H1a)}: More restrictive licenses are associated with higher total numbers of authors.
    \item \textbf{Hypothesis (H1b)}: More restrictive licenses (often ideologically motivated) are more likely to be retained.
\end{itemize}

\subsection{Social and Technical Choice}

Licenses, like other technologies and practices, spread through
communities.
Prior studies, such as those by \citet{ma2020methodology} and 
others~\cite{lamba2020heard, choice23}, used the
theory of social contagion~\cite{shoroye2015exploring} to estimate
the impact of exposure, susceptibility, and infectiousness on
developer choices regarding package (and tool) selection.
These studies found
that exposure measured by the number of overall deployments at the
time of choice, had a significant positive effect on 
selection.

By analogy, developers indifferent to or unaware of license
distinctions likely gravitate to licenses they have encountered
(prevalence-driven adoption), and compatibility with collaborators'
existing choices reinforces this (kinship-driven adoption).
\icsetrim{The choice of programming language also influences license
selection due to the culture and practices embedded within specific
developer communities.}
Programming language captures community norms and technical
compatibility: Python and JavaScript projects favor permissive licenses
like MIT or Apache~\cite{lerner2005economics}, while C/C++ projects
tend toward GPL~\cite{fitzgerald2006transformation}.
\icsetrim{The community size associated with a project can also impact
license decisions. \citet{tsay2014influence} found that projects with a
large number of forks are perceived as more valuable and trustworthy by
the community. Projects aiming to encourage collaboration and reuse to
maintain this momentum may favor permissive licenses. Similarly,
\citet{borges2018s} suggest that the number of stars signals community
approval, which, like forks, may influence the choice of more
permissive licenses.}
Community size (forks~\cite{tsay2014influence},
stars~\cite{borges2018s}) and the number of upstream/downstream
projects also plausibly influence license choice, with more interacting
projects favoring balanced (weak-copyleft or conditional) licenses over
extremes. Hence:
\begin{itemize}
    \item \textbf{Hypothesis (H2a)}: License choice is associated with the overall popularity of the license
    at adoption time.
    \item \textbf{Hypothesis (H2b)}: License choice is associated with programming language culture-specific norms.
    \item \textbf{Hypothesis (H2c)}: More permissive licenses are associated with larger community size
    (i.e., more upstream and downstream projects).
\end{itemize}

\subsection{Project Goals}

\icsetrim{While legal compliance sets essential boundaries, the
specific needs and objectives of a project further influence the
selection of an appropriate license. Developers often choose licenses
based on their preferences for the future use of the
software.}
Project-level objectives (e.g., preference for
widespread adoption versus strong copyleft
enforcement~\cite{sen2008determinants}) and the resulting trajectory
\cite{colazo2009impact} also matter.
We use four observable measures to operationalize trajectory:
commit count (activity~\cite{koch2002effort}),
file count (complexity/modularity~\cite{mockus2007large,bird2009putting}),
active-months (maturity~\cite{capiluppi2003characteristics}), and
burstiness (fluctuation, computed as active-months over lifetime).
\icsetrim{Projects with frequent commits tend to have more active
development cycles, making them more reliable and attractive for reuse.
Larger projects with more files are more likely to offer reusable
components, making license choice critical for balancing reuse and
intellectual property protection. Older, more mature projects are
often seen as more reliable, which could be influenced by the choice
of licenses that support long-term sustainability. Projects with high
burstiness may not focus heavily on legal concerns, and instead tend
to choose well-known, widely used licenses.}

\begin{itemize}
\item \textbf{Hypothesis (H3a)}: License choice is associated with project activity, approximated by number of commits.
\item \textbf{Hypothesis (H3b)}: More complex projects (higher number of files and blobs) are more likely to use restrictive licenses.
\item \textbf{Hypothesis (H3c)}: Restrictive licenses are associated with longer-term sustainability, approximated by project activity duration.
\item \textbf{Hypothesis (H3d)}: Burstiness, reflecting fluctuating activity, is associated with more permissive licenses.
\end{itemize}

\section{Comparisons to Prior Work}

Despite legal importance of license choice and the potential to stunt
innovation and poison supply chains, it appears that a significant
portion of OSS participants are ignoring the need to choose a license
or to verify license compatibility. 
We construct a preliminary
theory of how license choice might affect the project by reviewing relevant
literature, operationalizing potential affected factors,
and testing that theory. 
However, while our work is not the first to empirically study software licenses,
it differs from the literature in two important ways.

\textit{Comprehensive Identification of Licenses.}
Most prior research, such as that by \citet{wu2024large} and \citet{xu2023lidetector}, primarily depends on explicit license declarations found in metadata files.
Others, like \citet{feng2019open}, apply static analysis on binaries to detect embedded license texts.
However, these methods may overlook licenses that are not clearly stated or are stored in unconventional directories.
In contrast, we use a comprehensive dataset by \citet{jahanshahi2024license}, 
% which follows a more comprehensive approach 
compiled by scanning virtually the entire OSS ecosystem for any files containing ``license'' in their filepath.
Moreover, the authors use the winnowing algorithm, a robust method for matching license texts to known licenses, enhancing the accuracy of detecting both partial and full matches, even when the text is embedded or slightly modified.
This method captures not only standard license files but also other files potentially holding licensing details, ensuring that no relevant license information is missed.

\textit{Scale and Scope of Analysis.} 
Previous works often limit their scope to specific platforms (e.g., GitHub), a few package manager environments (e.g., NPM), or types of licenses (e.g., OSI-approved licenses), lacking a comprehensive, large-scale approach to detecting and analyzing licenses.
Our study expands this scope by analyzing essentially the entire open source landscape, ensuring a more comprehensive cross-platform understanding of licensing practices.

\icsetrim{In Table~\ref{tbl:litRev}, we compare our study's
methodology, coverage, scope, and scale with the most recent
comparable studies in the field.

\begin{table*}[t]
\centering
\caption{Latest Studies in Open Source Software Licensing}
\label{tbl:litRev}
\resizebox{0.95\textwidth}{!}{
\begin{tabular}{llllr}
    \toprule
    & \textbf{License Detection} & \textbf{Coverage} & \textbf{Scope} & \textbf{Scale} (Projects) \\
    \midrule
    \cite{wu2024large} & Package Manager API & License Files & Maven, NPM, PyPI, RubyGems, Cargo & 3,474,778 \\
    \cite{cui2023empirical} & NLP/Extract License Terms & All Files & GitHub - >1000 Stars & 16,341 \\
    \cite{wolter2023open} & GitHub API+NLP/Extract License Terms & All Files & GitHub - Cataloged by OpenHub & 1,000 \\
    \cite{xu2023lidetector} & NLP/Extract License Terms & All Files & GitHub - Popular Python Projects & 1,846 \\
    \multicolumn{5}{c}{\dotfill} \\
    \textbf{This Study} & \textbf{NLP/Matching Known Licenses} & \textbf{License Files} & \textbf{Entire OSS} & \textbf{131,171,379} \\
    \bottomrule
\end{tabular}}
\end{table*}}
Compared to recent licensing studies
(\citet{wu2024large}: 3.5M projects across 5 package managers;
\citet{cui2023empirical}: 16k high-star GitHub repos;
\citet{wolter2023open}: 1k OpenHub-catalogued repos;
\citet{xu2023lidetector}: 1.8k popular Python repos), our analysis
covers $131$M deforked WoC projects -- two orders of magnitude
larger and not filtered by platform, stars, or package-manager
presence.
